\subsection{Widerstandsberechnung für invertierenden Verstärker}
% Alt: Aufgabe 3
\Aufgabe{
    \label{Aufg3}
    Gegeben ist die abgebildete Schaltung mit einem invertierenden
    Verstärker. Dabei ist der Widerstand $R_\mathrm{2} = 1\,\mathrm{k \Omega}$. Was für Widerstandswerte
    muss $R_\mathrm{1}$ aufweisen, damit die Verstärkung der gegebenen Schaltung -10, -50 und -100 beträgt?

    \begin{figure}[H]
        \centering
        \input{Tikz/src/FigWiderstandIV.tex}\alttext{Das Bild zeigt ein elektrisches Schaltbild mit einem Operationsverstärker.Links ist eine Eingangsspannung \(U_\mathrm{E}\) mit einem Pfeil nach unten eingezeichnet.Diese Eingangsspannung ist über den Widerstand \(R_1\) mit dem Minus-Eingang (invertierender Eingang) verbunden. Vom Knoten zwischen \(R_1\) und dem invertierenden Eingang führt ein weiterer Zweig nach oben zu einem zweiten Widerstand \(R_2\).Dieser Widerstand verbindet den Ausgang des Operationsverstärkers mit dem invertierenden Eingang. Dadurch entsteht eine Rückkopplung. Der Puls-Eingang (Nichtinvertierende Eingang) des OPV ist nach unten direkt mit Masse verbunden. Rechts befindet sich der Ausgang des Operationsverstärkers. Dort ist die Ausgangsspannung \(U_\mathrm{A}\) mit einem Pfeil nach unten eingezeichnet.}
        \label{fig:FigWiderstandIV}
    \end{figure}
}

\Loesung{
    Gemäß Formel invertierender Verstärker:
    \begin{align*}
        V = - \frac{R_2}{R_1}
    \end{align*}
    Formel umstellen nach $R_\mathrm{1}$:
    \begin{align*}
        R_1 = - \frac{R_2}{V}
    \end{align*}
    \begin{itemize}
        \item $V = - 10$
              \begin{align*}
                  R_1 =\mathrm{- \frac{1\,k\Omega}{- 10} = 100\,\Omega}
              \end{align*}
        \item $V = - 50$
              \begin{align*}
                  R_1 =\mathrm{- \frac{1\,k\Omega}{- 50} = 20\,\Omega}
              \end{align*}
        \item $V = - 100$
              \begin{align*}
                  R_1 =\mathrm{- \frac{1\,k\Omega}{- 100} = 10\,\Omega}
              \end{align*}
    \end{itemize}

    Siehe: Abschnitt \ref{fig: Verschaltung Verstaerker} und Tabelle \ref{tab:Grundschaltungen}
}